\documentclass[11pt]{article}
\usepackage[margin=2cm]{geometry}
\usepackage{amsmath,amssymb}
\usepackage{graphicx}
\usepackage{caption}
\usepackage{booktabs}
\usepackage[dvipsnames,table]{xcolor}
\usepackage{titlesec}
\usepackage{enumitem}
\usepackage{parskip}
\usepackage{mdframed}
\usepackage{hyperref}

\definecolor{deep}{HTML}{065A82}
\definecolor{teal}{HTML}{1C7293}
\definecolor{navy}{HTML}{21295C}
\definecolor{ice}{HTML}{EAF1FB}
\definecolor{ink}{HTML}{1A2230}

\hypersetup{colorlinks=true, linkcolor=deep, urlcolor=deep}
\titleformat{\section}{\large\bfseries\color{deep}}{\thesection.}{0.5em}{}
\titleformat{\subsection}{\normalsize\bfseries\color{teal}}{}{0em}{}
\setlist[itemize]{leftmargin=1.4em, itemsep=2pt, topsep=2pt}

\graphicspath{{D:/DINOSR/Claude/PaperRun_claude/dino_sr_contrastive/docs/explainer/figs/}}

\newmdenv[backgroundcolor=ice, linecolor=deep, linewidth=0.8pt,
          roundcorner=4pt, skipabove=6pt, skipbelow=6pt,
          innertopmargin=6pt, innerbottommargin=6pt]{callout}

\begin{document}
\pagestyle{empty}

\begin{center}
{\LARGE\bfseries\color{navy} What is the model actually looking at?}\\[3pt]
{\large\color{teal} Making the self-supervised 4D-STEM ``grain map'' explainable
in physical terms}\\[2pt]
{\small\color{ink} A non-technical guide to the interpretability tests
$\cdot$ IMC \& NaPHI $\cdot$ Daniel Khaykelson}
\end{center}
\vspace{2pt}
\hrule height 1pt
\vspace{8pt}

\section{The setup, and the problem}
Our model (a self-supervised network called DINO) looks at the electron
diffraction pattern recorded at every probe position in a 4D-STEM scan and
sorts each position into one of $\sim$14 groups---producing a ``grain map''
\emph{with no labels and no chemistry input}. It works, but it is a black box:
we do not know \emph{which physical property} defines the groups
(crystallinity? thickness? crystal orientation? something new?). The four
tests below answer that, each from a different angle, so that together they
triangulate a trustworthy physical explanation.

\section{The four tests (plain idea + what it measures)}

\subsection{1. Probing --- what the model pays attention to}
\textbf{Plain idea:} if we can ``read off'' a known property (say
crystallinity) from the model's internal notes for each point, then the model
is keeping track of it.\\
\textbf{What we measure:} we fit a simple linear predictor from the model's
128-number fingerprint to each measured property; the score $R^2$ (0--1) says
how strongly that property is encoded.

\subsection{2. Ablation --- what the model depends on}
\textbf{Plain idea:} like covering part of a photo to see whether you can still
recognise it.\\
\textbf{What we measure:} we blank out one ingredient of the diffraction
pattern (overall brightness, fine angular detail, or inner vs.\ outer rings),
re-run the model, and compare the new map to the original (agreement = ARI).
If the map survives the edit, that ingredient was not needed; if it falls
apart, the model depended on it.

\subsection{3. Classical check --- is the map actually new?}
\textbf{Plain idea:} does the AI just redraw what a standard, decades-old
analysis already shows, or something genuinely new?\\
\textbf{What we measure:} we cluster classical 4D-STEM features (a virtual
dark-field image, the azimuthally-averaged profile, and PCA/NMF
decompositions of the patterns) into the same number of groups, and measure
how much they agree with the AI map (ARI/AMI).

\subsection{4. Orientation check --- is it just sorting by crystal direction?}
\textbf{Plain idea:} a tilted crystal diffracts differently; maybe the model
is only sorting by grain \emph{orientation}.\\
\textbf{What we measure:} where a crystal structure (CIF) is available, we
index the orientation at each point (ACOM) and measure agreement with the AI
map. We also use Grad-CAM to see \emph{which part} of the pattern the model
looks at.

\begin{callout}
\textbf{How to read the two numbers.}
\;$R^2$ --- how predictable a property is from the model
(\,0 = ignored\, $\rightarrow$\, 1 = fully encoded\,).\quad
\textbf{ARI / AMI} --- how much two maps agree, corrected for chance
(\,0 = no better than random\, $\rightarrow$\, 1 = identical maps\,).
\end{callout}

\section{What we found}
\begin{itemize}
\item The groups are driven by the \textbf{overall scattered (diffracted)
intensity} together with the \textbf{inner, low-angle (low-$q$) 2-D
diffraction structure} --- the arrangement of the inner rings/spots
($q \lesssim 2.1$\,nm$^{-1}$, i.e.\ $d \gtrsim 4.7$\,\AA).
\item It is \textbf{not} based on crystal orientation: agreement with the
indexed zone axis is essentially zero (AMI $\le 0.06$), because the model is
trained to be rotation-invariant by design.
\item A 1-D radial profile alone is \textbf{not enough} --- averaging away the
angles collapses the map; the full 2-D pattern is needed.
\item \textbf{Uniqueness is material-dependent.} For IMC, no classical method
reproduces the map (best agreement ARI $\approx 0.15$): the AI map is
\emph{distinctive}. For NaPHI, classical methods get about halfway
(ARI $\approx 0.5$): there the AI stays closer to a classical decomposition.
\item Logarithmic intensity stretching (an image-prep choice) keeps the same
physical basis but \emph{sharpens} it --- more, cleaner groups.
\end{itemize}

\section{What it means (and one caveat)}
The model groups regions by \emph{how much, and what, they diffract at low
angles} --- crystallinity, molecular packing, and sample thickness. That is a
physically sensible axis, not an arbitrary black box. For IMC it is genuinely
distinctive (a standard virtual-detector or PCA map does not reproduce it);
for NaPHI it stays closer to classical analysis.

\textbf{Caveat:} ``scattered intensity'' blends crystallinity with sample
thickness --- both raise the total diffracted signal. An independent
thickness map (e.g.\ a $t/\lambda$ estimate) would be needed to separate the
two. Orientation/phase analysis is folded in automatically wherever a CIF is
available; the conclusions above hold with or without it.

\begin{figure}[h]
\centering
\includegraphics[width=\linewidth]{cross_sample.png}
\caption*{\small Cross-sample summary. \textbf{(a)} which properties are
encoded (probe $R^2$); \textbf{(b)} causal ablations (low bar = the model
depends on that signal); \textbf{(c)} how well classical methods reproduce the
AI map (all $\le 0.18$ for IMC). Bars: SI4, SI3 (linear), SI3 (log-stretch).}
\end{figure}

\begin{figure}[h]
\centering
\includegraphics[width=0.9\linewidth]{imc_dino_vs_classical.png}
\caption*{\small IMC: the AI map (far left) vs.\ classical clusterings.
Colours are arbitrary --- what matters is the spatial shape. Classical
methods catch the coarse regions but miss the fine grain structure the AI
resolves.}
\end{figure}

\end{document}
